\documentclass[11pt]{article}

\usepackage[T1]{fontenc}
\usepackage[utf8]{inputenc}
\usepackage{lmodern}
\usepackage{microtype}
\usepackage{amsmath,amssymb,amsthm,mathtools}
\usepackage{enumitem}
\usepackage{xcolor}
\usepackage{hyperref}
\usepackage[nameinlink,capitalise,noabbrev]{cleveref}

\emergencystretch=6em

\hypersetup{
  colorlinks=true,
  linkcolor=blue!45!black,
  citecolor=blue!45!black,
  urlcolor=blue!45!black
}

\newtheorem{theorem}{Theorem}[section]
\newtheorem{lemma}[theorem]{Lemma}
\newtheorem{proposition}[theorem]{Proposition}
\newtheorem{corollary}[theorem]{Corollary}
\newtheorem{definition}[theorem]{Definition}
\newtheorem{assumption}[theorem]{Assumption}
\newtheorem{remark}[theorem]{Remark}

\newcommand{\HE}{Homotopical Encryption}
\newcommand{\HCP}{\mathrm{HCP}}
\newcommand{\CVP}{\mathrm{CVP}}
\newcommand{\MCA}{\mathrm{MCA}}
\newcommand{\RS}{\mathrm{RS}}
\newcommand{\F}{\mathbb{F}}
\newcommand{\R}{\mathbb{R}}
\newcommand{\Z}{\mathbb{Z}}
\newcommand{\Gr}{\operatorname{Gr}}
\newcommand{\Path}{\operatorname{Path}}
\newcommand{\Fib}{\operatorname{Fib}}
\newcommand{\image}{\operatorname{im}}
\newcommand{\coker}{\operatorname{coker}}
\newcommand{\dist}{\operatorname{dist}}
\newcommand{\Public}{\mathsf{Public}}
\newcommand{\Hidden}{\mathsf{Hidden}}
\newcommand{\Lean}{\textsc{Lean}}
\newcommand{\lean}[1]{\begingroup\Urlmuskip=0mu plus 1mu\path{#1}\endgroup}
\newenvironment{leanledger}{\small\sloppy}{\par}
\newcommand{\todo}[1]{\textcolor{red!65!black}{[TODO: #1]}}


\title{Homotopical Encrypt: Hidden Hodge Operators and the Harmonic Cochain Problem}
\author{Chloe van der Vlugt}
\date{}

\begin{document}

\maketitle

\begin{abstract}
We introduce the Harmonic Cochain Problem (HCP), the problem of recovering the
harmonic cochain \(h\in\ker(\Delta)\) selected by a hidden Hodge--Laplacian
\(\Delta\) from a public homotopy class, without determining \(\Delta\) or
\(\ker(\Delta)\). A hidden representative determines restriction data,
adjacent coboundaries \(\delta_0,\delta_1\), the degree-one Hodge--Laplacian
\(\Delta_1=\delta_0\delta_0^\ast+\delta_1^\ast\delta_1\), and the Hodge
projection \(\pi_H:C^1\to\ker(\Delta_1)\) used for recovery; the public class
records only a classifier shadow, while the observed cochain has the form
\(b = h + e\), where
\(e \in \image(\delta_0) \oplus \image(\delta_1^\ast)\) is Hodge-structured noise.
For every explicitly generated lattice
\(\Lambda_B=\image(B)\), there is a linear operator \(\Delta\) with
\(\ker(\Delta)=\Lambda_B\); hence a presented-operator closest-harmonic solver
returns a closest vector for the CVP instance \((B,t)\). Conversely, a correct
known-lattice CVP oracle does not uniformly solve public-data HCP when publicly
equivalent hidden operators may disagree on harmonicity. We isolate the
semantic and axiomatic surfaces needed
for this statement, including fibrational identity data, homotopy classes,
exact invariants, spectral/Tsirelson admissibility, and Heyting-style oracle
judgments. We then define the public artifact stack of cochain commitment
payloads, commitment axioms, proof-carrying artifacts, and lossy projections,
and construct the finite projection endpoint as a Reed--Solomon
mutual-correlated-agreement statement. The definitions, theorem boundaries, and
Lean correspondence are formalized and explicitly indexed in the appendix
ledger.
\end{abstract}

\tableofcontents

\section{Introduction}

\subsection{Motivation}

\HE{} is organized around the following security shape:
\[
\begin{gathered}
\text{public handles}\\
\text{and verifier-facing artifacts}
\end{gathered}
\quad\not\Rightarrow\quad
\text{the operative Hodge--Laplacian kernel}.
\]
The private recovery object is the harmonic cochain selected by a hidden Hodge
operator, instead of a scalar or a vector secret in a published lattice or algebraic object. 
Public verification therefore needs finite artifacts, while security
must not collapse to recovering the private representative or the hidden
operator.

\subsection{Contributions}

This paper develops the following contributions.
\begin{enumerate}[label=(\roman*)]
  \item The Harmonic Cochain Problem (HCP): recovery of the harmonic cochain
    selected by a hidden Hodge operator.
  \item A Hodge-structured noise model with cochain decomposition
    \(b = h + e\), where \(h \in \ker(\Delta)\) and
    \(e \in \image(\delta) \oplus \image(\delta^\ast)\).
  \item A CVP comparison in which an explicit lattice is realized as a kernel
    \(\ker(\Delta)\), a harmonic solver yields a lattice vector, and the
    closest-harmonic contract recovers the closest-vector predicate.
  \item A hidden-operator ambiguity theorem: publicly equivalent Hodge operators
    may disagree on harmonicity, so a correct explicit-lattice CVP oracle does
    not uniformly solve public-data HCP.
  \item A Lean formalization of the definitions, reductions, projection chain,
    gates, and theorem boundaries.
  \item A projection chain through Plucker shadows and cochain commitments to a
    canonical Reed--Solomon endpoint with an MCA bad-challenge bound.
\end{enumerate}

\section{The Harmonic Cochain Problem}

This section states the central search problem and the hidden-operator model in
which it is used.

\begin{definition}[Harmonic Cochain Problem]\label{def:hcp}
  The \emph{Harmonic Cochain Problem} (HCP) asks, given public
  cochain/class-level data associated to a Hodge system, to recover the
  harmonic cochain determined by the operative Laplacian kernel.
\end{definition}

In the \HE{} public-sample model, the visible cochain has the form
\[
  b = h + e,
  \qquad
  h \in \ker(\Delta),
  \qquad
  e \in \image(\delta) \oplus \image(\delta^\ast).
\]
Here \(h\) is the harmonic cochain, equivalently the global-section component,
and \(e\) is Hodge-structured noise. The hidden-operator condition is that the
public data does not determine a unique \(\Delta\) or \(\ker(\Delta)\).

\subsection{Public And Hidden Data}

\begin{definition}[Public observation]\label{def:public-observation}
  A public observation is the data exposed to a verifier or adversary before
  choosing a hidden representative. In the abstract model it consists of a
  public class or shadow value \(p\), together with cochain-facing data \(b\).
  In the \HE{} instantiation, this public layer is represented by public
  shadows, cochain commitments, topology hashes, oracle judgments, and
  proof-carrying artifacts.
\end{definition}

The public observation is intentionally not the hidden Hodge system. It is a
classifier of a public equivalence class of possible hidden representatives.
This is the point at which the fibrational and homotopical semantics enter the
model: public equality is equality after applying a public classifier, not
strict equality of hidden representatives. Section~\ref{sec:mathematical-framework}
develops this semantic layer.

\begin{definition}[Hidden Hodge witness]\label{def:hidden-hodge-witness}
  A hidden Hodge witness for a public observation \((p,b)\) consists of
  hidden operator data \(\Delta\), a harmonic cochain \(h\), and
  Hodge-structured noise \(e\), such that
  \[
    \Delta \in \Hidden(p),
    \qquad
    h \in \ker(\Delta),
    \qquad
    e \in \image(\delta) \oplus \image(\delta^\ast),
    \qquad
    b = h + e.
  \]
  Here \(\Hidden(p)\) denotes the fiber of hidden Hodge--Laplacian
  representatives compatible with the public class value \(p\).
\end{definition}

Concretely, the hidden data includes the fibration representative, restriction
maps, the induced Laplacian, its harmonic kernel, and private witness/cochain
data. The public data may bind to this hidden data through commitments or
proof-carrying artifacts, but it does not present the operative kernel itself.

\begin{definition}[HCP solution]\label{def:hcp-solution}
  A solution to an HCP instance \((p,b)\) is a cochain \(h'\) that is the
  harmonic component selected by the operative hidden witness. In the exact
  recovery form used by the hidden-operator separation, \(h' = h\) for every
  hidden witness \((\Delta,h,e)\) compatible with the public observation.
\end{definition}

The exact recovery formulation is deliberately strong. It isolates the
ambiguity proved in the Lean development: the same public cochain may be
compatible with hidden operators that disagree about harmonicity. Thus a solver
that sees only public data cannot be uniformly sound unless the public
observation determines enough hidden-operator structure.

\subsection{Recovery Target}

The recovery target is the harmonic cochain selected by the operative kernel,
not necessarily the entire kernel subspace. This distinction keeps HCP parallel
to CVP: CVP asks for a closest vector in a presented lattice, while HCP asks for
a harmonic cochain in the operative Hodge kernel.

In the metric variant used for comparison with CVP, one may ask for the closest
harmonic cochain to \(b\). Once \(\Delta\) or a basis for \(\ker(\Delta)\) is
explicitly presented, this becomes a closest-vector question on the harmonic
lattice. The hidden-operator model withholds precisely that presentation.

\subsection{Computational Model}

The computational model separates three layers.
\begin{enumerate}[label=(\roman*)]
  \item The mathematical layer defines public projections, hidden operator
    fibers, harmonicity, and kernel ambiguity.
  \item External hardness assumptions and cited coding-theory theorems enter at
    explicitly labeled theorem boundaries.
  \item Runtime obligations such as byte parsing, canonical encodings, hash
    domain separation, proof verification, and state transition are outside the
    mathematical HCP definition.
\end{enumerate}

This paper focuses on the first two layers. The runtime layer is part of the
protocol and proof-carrying-verification story developed in the sequel.

\section{CVP Comparison And Kernel Ambiguity}

CVP asks for a closest vector in an explicitly presented lattice. HCP asks for
a harmonic cochain in the operative Hodge kernel, and in the \HE{} model that
kernel is not explicitly presented by public data.

\subsection{Known-Lattice CVP Surface}

An explicit-lattice CVP instance consists of a generator matrix \(B\) and a
target vector \(t\). The presented lattice is
\[
  \Lambda_B=\{Bx:x\in R^n\}.
\]
A correct CVP answer is a vector \(v\in\Lambda_B\) such that
\[
  \dist(t,v)\le \dist(t,w)
  \qquad
  \text{for every }w\in\Lambda_B.
\]
This is the separation between the known-lattice interface and the
hidden-operator interface: a CVP oracle receives \(B\) and \(t\), while a
public HCP procedure sees only the declared handles and verifier-facing
artifacts unless extra hidden-operator information is supplied.

\subsection{Reduction Direction}

The formalization separates three ingredients. The first is pure linear
algebra: every presented lattice \(\Lambda_B=\image(B)\) can be realized as
the kernel of some matrix \(\Delta\). The second is the presented-operator
Hodge contract: if a solver is sound for the displayed operator \(\Delta\),
then its output is harmonic for that \(\Delta\), hence lies in
\(\ker(\Delta)=\Lambda_B\). The third is the closestness contract owned by the
CVP interface: only after adding minimization among harmonic vectors does the
membership statement become a closest-vector statement. Lean keeps these
surfaces separate.

\begin{proposition}[Kernel realization for explicit lattices]\label{prop:hcp-to-cvp}
  Let \(B\) be an explicit generator matrix over a field and set
  \(\Lambda_B=\image(B)\). There exists a matrix \(\Delta\) such that
  \[
    \ker(\Delta)=\Lambda_B .
  \]
  Consequently, any solver whose output is \(\Delta\)-harmonic for every
  presented \(\Delta\) returns a vector in the explicit lattice
  \(\Lambda_B\) when run on this realized operator. If that solver is also
  closest among \(\Delta\)-harmonic vectors, then its output is a CVP solution
  for \((B,t)\).
\end{proposition}

\begin{proof}
  Let \(B\) be an explicit generator matrix for a lattice
  \[
    \Lambda_B=\{Bx:x\in R^n\}\subseteq R^n
  \]
  and let \(t\in R^n\) be the CVP target. We first show the following
  linear-algebra fact: there is a matrix \(\Delta\) with
  \[
    \ker(\Delta)=\image(B)=\Lambda_B .
  \]
  Put \(V=\image(B)\). Since \(R^n\) is finite-dimensional over a field, choose
  a complementary subspace \(W\) with
  \[
    V+W=R^n,
    \qquad
    V\cap W=0 .
  \]
  Thus every \(y\in R^n\) has a unique decomposition
  \[
    y=v+w,
    \qquad
    v\in V,\quad w\in W.
  \]
  Equivalently, there is a unique element \(p(y)\in W\) such that
  \(y-p(y)\in V\). Define
  \[
    P(y)=p(y).
  \]
  The uniqueness of the \(V\oplus W\) decomposition gives linearity: for
  \(y,z\in R^n\), the vector \(P(y)+P(z)\) lies in \(W\), and
  \[
    (y+z)-(P(y)+P(z))
    =
    (y-P(y))+(z-P(z))\in V,
  \]
  so \(P(y+z)=P(y)+P(z)\). Similarly, for \(a\in R\),
  \(aP(y)\in W\) and
  \[
    ay-aP(y)=a(y-P(y))\in V,
  \]
  so \(P(ay)=aP(y)\).

  We now compute the kernel. If \(y\in V\), then \(0\in W\) and \(y-0\in V\),
  hence uniqueness gives \(P(y)=0\). Conversely, if \(P(y)=0\), then
  \(y-P(y)=y\in V\). Therefore
  \[
    \ker P=V=\image(B).
  \]
  Let \(\Delta\) be the matrix of the linear map \(P\). Then, for every
  \(y\in R^n\),
  \[
    \Delta y=0
    \quad\Longleftrightarrow\quad
    y\in\ker P
    \quad\Longleftrightarrow\quad
    y\in\image(B)
    \quad\Longleftrightarrow\quad
    \exists x,\ Bx=y .
  \]
  This is the content of the Lean lemma
  \lean{exists_matrix_ker_eq_range}.

  Now apply a presented-operator Hodge solver to \(t\) under this realized
  operator. In Lean the solver is a function \(S:R^n\to R^n\), and the
  presented-operator soundness hypothesis is
  \[
    \forall \Delta\, b,\qquad
    \Delta(S(b))=0.
  \]
  Applying this with the \(\Delta\) constructed above and \(b=t\) gives
  \[
    \Delta(S(t))=0.
  \]
  By the kernel calculation, there exists \(x\) with
  \[
    Bx=S(t).
  \]
  Hence \(S(t)\in\Lambda_B\). This is exactly the theorem
  \lean{Hodge_Hardness_Reduction_Proof}: a solver that is harmonic for every
  presented operator yields a lattice vector for every explicit-lattice
  instance once the lattice has been realized as a kernel.

  Finally, suppose in addition that the solver satisfies the closest-harmonic
  contract
  \[
    \forall h',\quad
    \Delta h'=0
    \Longrightarrow
    \dist(t,S(t))\le \dist(t,h').
  \]
  Since the \(\Delta\)-harmonic vectors are exactly the vectors of
  \(\Lambda_B\), this inequality is precisely
  \[
    \dist(t,S(t))\le \dist(t,w)
    \qquad
    \text{for every }w\in\Lambda_B.
  \]
  Together with \(S(t)\in\Lambda_B\), this is the CVP correctness predicate
  \lean{IsClosestLatticeVector} for the explicit instance \((B,t)\).
\end{proof}

Once the hidden operator data has been fixed or presented, the closest harmonic
cochain problem is precisely a CVP-style problem on \(\ker(\Delta)\). Thus CVP
is a known-lattice shadow of the Hodge recovery problem.

\begin{proposition}[Projection-model Hodge recovery is CVP-shaped]\label{prop:hodge-projection-is-cvp}
  Fix a hidden harmonic subspace \(W\subseteq R^n\) with orthogonal projection
  \(P_W\), and let \(\Delta_W=I-P_W\). A cochain \(h\) is
  \(\Delta_W\)-harmonic exactly when \(h\in W\). Therefore, if \(h\) is
  harmonic and closest to \(b\) among all harmonic cochains, then \(h\) is the
  CVP solution for the target \(b\) in the lattice/subspace \(W\).
\end{proposition}

\begin{proof}
  In the projection model,
  \[
    \Delta_W b=(I-P_W)b,
  \]
  equivalently the projection of \(b\) to \(W^\perp\). Hence
  \[
    \Delta_W h=0
    \quad\Longleftrightarrow\quad
    h\in W.
  \]
  The remaining statement is definitional: the closest harmonic representative
  is precisely a vector that is \(\Delta\)-harmonic and minimizes
  \(\dist(b,-)\) among all \(\Delta\)-harmonic vectors.
\end{proof}

\subsection{Oracle Insufficiency}

The strictness direction uses a different contract. A public Hodge solver sees
the cochain-facing observation and may call an explicit-lattice CVP oracle, but
it is not given the hidden operator. Soundness asks that the solver's output be
harmonic for every operative hidden \(\Delta\). Exactness on already-harmonic
inputs asks that, if the observation \(b\) is already harmonic for the operative
\(\Delta\), the solver return \(b\) unchanged.

\begin{proposition}[Known-lattice CVP oracle insufficiency]\label{prop:cvp-insufficient}
  In the hidden-operator model, an oracle for explicit known-lattice CVP does
  not by itself determine the operative Hodge--Laplacian kernel. Consequently,
  such an oracle does not uniformly solve HCP from public data alone.
\end{proposition}

\begin{proof}
  Assume \(\dim R^n\ge 2\), and choose an index \(i\). Let
  \[
    b=e_i,\qquad
    \Delta_1=0,\qquad
    \Delta_2=\operatorname{diag}(1,0,\ldots,0),
  \]
  where the nonzero diagonal entry of \(\Delta_2\) is placed at \(i\). Then
  \[
    \Delta_1 b=0,
    \qquad
    \Delta_2 b=b\ne 0.
  \]
  Thus \(b\) is harmonic for \(\Delta_1\) but not harmonic for
  \(\Delta_2\). If the public projection is lossy enough to identify
  \(\Delta_1\) and \(\Delta_2\), the same cochain-facing observation \(b\) is compatible
  with two hidden Hodge operators that disagree about whether \(b\) lies in the
  harmonic kernel.

  Now suppose, for contradiction, that there were a public-data-only Hodge
  solver \(S\) which may call a correct explicit-lattice CVP oracle \(O\), and
  which is uniformly sound and exact on already-harmonic inputs. Since \(b\) is
  harmonic for \(\Delta_1\), exactness forces
  \[
    S(O,b)=b.
  \]
  But soundness for the hidden operator \(\Delta_2\) forces
  \[
    \Delta_2 S(O,b)=0.
  \]
  Combining the two equalities gives \(\Delta_2 b=0\), contradicting the
  construction of \(b\). Hence no such public-data-only solver exists, even
  assuming the availability of a correct oracle for explicitly presented CVP
  instances.
\end{proof}

The obstruction is kernel ambiguity: the same public observation may be
compatible with multiple hidden operators whose harmonic kernels differ. A CVP
oracle can solve a closest-vector problem after the lattice is presented; it
does not select which hidden kernel public data was meant to present.

\begin{proposition}[Concrete public ambiguity]\label{prop:constant-public-ambiguity}
  In every dimension at least two, the constant public projection is
  public-compatible Hodge ambiguous, hence kernel-ambiguous and lossy.
\end{proposition}

\begin{proof}
  Use the same witnesses as in the proof of
  \cref{prop:cvp-insufficient}: \(b=e_i\), \(\Delta_1=0\), and
  \(\Delta_2=\operatorname{diag}(1,0,\ldots,0)\). The constant projection sends
  every hidden operator to the same public value, so \(\Delta_1\) and
  \(\Delta_2\) are publicly equivalent. Since \(b\) is harmonic for
  \(\Delta_1\) and not for \(\Delta_2\), this gives public-compatible Hodge
  ambiguity. The implications of \cref{prop:ambiguity-hierarchy} then give
  kernel ambiguity and lossiness.
\end{proof}

\subsection{Structural Comparison}

The comparison separates three surfaces:
\begin{enumerate}[label=(\roman*)]
  \item known-lattice closest-vector search;
  \item hidden-operator HCP;
  \item public projection/proximity verification.
\end{enumerate}

The first two surfaces are the hardness story. The third is a verifier-facing
theorem: given accepted hidden/geometric data, later sections construct finite
public shadows and an RS/MCA endpoint without claiming that the projection is
injective.

\section{Axioms, Semantics, And Public Data}\label{sec:mathematical-framework}

This section fixes the semantic and axiomatic surfaces used by the problem.
It comes before the public artifact and commitment layer because those objects
are only meaningful after the public/private boundary is fixed.

\subsection{Fibrational Semantics}

The semantic layer is motivated by Helfer's first-order homotopical logic
\cite{helfer2019fol}. In this paper, the corresponding Lean files are used as
a formalized interface for the semantic structure needed by \HE{}, not as a
claim that every external theorem from Helfer's paper has been reproved.

\subsection{Fibrational Identity Objects}

Public equality is not strict equality of hidden representatives. It is mediated
by public classifiers, path data, and equivalence of the observable structure.
This is the semantic source of the hidden-operator model used in HCP.

\subsection{Homotopy Classes}

Public shadows may identify multiple hidden representatives while preserving
the invariants required by a verifier. At this stage the shadow is only a
public classifier or fingerprint; finite-code projection is introduced later.

\subsection{Axiom Surfaces}

The Lean development separates assumptions and theorem surfaces into small
modules. The relevant early interfaces are:
\begin{enumerate}[label=(\roman*)]
  \item foundational cochain and Hodge assumptions;
  \item fibration and homotopy-class interfaces;
  \item Heyting/oracle semantics for public judgments;
  \item sheaf-coherence and spectral interfaces;
  \item cochain-commitment payload and security assumptions.
\end{enumerate}

\subsection{Exact Structural Invariants}

The exact gates used by the Lean development are not a separate source of
structure. They package the same foundational and sheaf-coherence hypotheses
at the point where an object is admitted. In Lean, \texttt{ExactGateBundle} is
the foundation interface and \texttt{SheafExactGateBundle} is the
sheaf-coherence interface. The theorem
\texttt{structural\_preconditions\_of\_sheaf\_exact\_gates} extracts the two
structural consequences used later: the cochain differential squares to zero
and the Hodge Laplacian is the one induced by the coboundary data.

\begin{proposition}[Exact structural preconditions]\label{prop:exact-structural-preconditions}
  An accepted sheaf-coherence surface supplies
  \[
    d^2=0
    \qquad\text{and}\qquad
    \Delta=d^\ast d+dd^\ast
  \]
  for the cochain complex under discussion.
\end{proposition}

\subsection{Heyting And Oracle Judgment Surface}

The public judgment layer is also fixed here, because HCP uses public
classifiers and oracle-visible predicates before any finite-code projection is
introduced. Lean separates an \texttt{OracleQuery} from transaction cochain
commitments: its \texttt{id\_commit} is an identity hash obtained from the
public Hodge projection and quantization path, not a cochain commitment
payload.

The oracle predicate checks public compatibility data such as the vertex
identifier, Betti number, identity hash, and Tsirelson-style spectral judgment.
The induced judgment is then mapped into the three-valued Heyting surface used
by the acceptance predicate. The stability axiom \texttt{OracleEta2} records
that equal public identity inputs produce the same identity hash; it is a
determinism condition for public identity observations, not a uniqueness or
binding theorem for cochain commitments.

This modular split is part of the mathematical statement. It prevents runtime
facts, cryptographic assumptions, public classifiers, and Hodge-theoretic
claims from being conflated. Exact and oracle invariants are used here as
mathematical acceptance surfaces. Consensus composition and chain/runtime
enforcement are recorded only at the formalization boundary; their full
CS/protocol treatment is deferred to the sequel.

\section{Public Artifact Stack And Lossy Projection}

Commitments and public shadows are the first concrete public-facing objects in
the paper. They explain what public data can carry without revealing the hidden
Hodge representative.
They are the scaffolding for transparency in action and anonymity in identity:
actions are exposed through auditable public artifacts, while identities remain
inside hidden representatives and public equivalence classes.

\subsection{Public, Artifact, And Private Layers}

The public-facing stack has three mathematical levels. The fully public handle
level contains hashes, roots, topology hashes, nullifiers, public summaries, and
oracle judgments. The verifier-facing artifact level contains public shadows,
cochain commitments, proof-carrying cochain payloads, and certificate data. The
private witness level contains the fibration representative,
transport/restriction data, full cochains and openings, and the Hodge operator
whose kernel selects the harmonic component. The paper's claims are stated so
that public handles and verifier-facing artifacts bind to hidden witnesses
without turning the hidden witness into public data. This is compatible with the
sheaf-theoretic attacker-observation viewpoint of
\cite{inoue2026cryptographicSheaves}; the associated security games and
runtime observation model belong to the subsequent protocol paper. The present
paper uses only the mathematical public/private separation described below.

\subsection{The Artifact Stack}

The public artifact stack has the following shape:
\[
\begin{aligned}
  &\text{hidden/private object} \\
  &\quad\to \text{public shadow or compressed public datum} \\
  &\quad\to \text{commitment/root/hash artifact} \\
  &\quad\to \text{proof-carrying artifact} \\
  &\quad\to \text{verifier judgment}.
\end{aligned}
\]
Each layer binds to or summarizes the preceding layer. A hash does not replace
the mathematical object, a commitment does not replace the witness, and a proof
artifact does not replace the proof obligation. The purpose of the stack is to
let public verification refer to hidden sheaf/cochain data without exposing it.
In this sense, the stack separates what must be publicly accountable from what
must remain representative-private.

Hashes serve as public names, integrity bindings, and fixed-size pointers for
larger artifacts. A topology hash names the public combinatorial support; a
root commits to a batch or storage view; a transcript hash domain-separates a
proof statement. None of these hashes asserts validity by itself. Validity
comes from the predicate or proof attached to the hashed object.

\subsection{Cochain Commitments}

Cochain commitments provide the public handle on hidden cochain data. The Lean
surface treats the canonical payload as a mathematical object with a modulus,
edge values, declared edge count, and topology hash.

\subsection{Commitment Map}

\begin{definition}[Cochain commitment payload]\label{def:cochain-commitment-payload}
  A cochain commitment payload is a tuple
  \[
    (q,n,c,\tau)
  \]
  consisting of a modulus \(q\), a declared edge count \(n\), a list \(c\) of
  integer edge values, and a topology hash \(\tau\). In the canonical Lean
  payload this is \texttt{CochainCommitmentData}.
\end{definition}

The intended commitment map is the transaction-facing coboundary map
\[
  (m,r) \longmapsto \delta(m\Vert r) \in C^1 \pmod q,
\]
stored together with the topology hash of the supporting complex. The payload
therefore preserves the accepted modulus, the edge-vector shape, the public
topology binding, and the serialized \(C^1\)-level commitment values. It does
not expose the hidden fibration representative, private transport data, the
operative Hodge--Laplacian, or the harmonic kernel.

\begin{definition}[Commitment well-formedness]\label{def:commitment-well-formedness}
  A cochain commitment payload is well formed when its modulus is the supported
  canonical modulus, its declared edge count agrees with the actual edge-value
  list length, and the edge list is nonempty.
\end{definition}

The security interface for this commitment layer consists of binding, hiding,
homomorphic addition, homomorphic scalar multiplication, topology binding, and
partial-opening soundness. The blinded-commitment, vector-commitment, and
cohomological-MAC interfaces are separate axiom records rather than consequences
silently bundled into the payload shape.

\subsection{Commitment And Artifact Invariants}

The commitment-facing invariants attach to the commitment payload and to the
public artifact stack, not to the proximity projection. Payload
well-formedness is a shape condition: it asserts the supported modulus,
agreement between the declared edge count and the edge-value list, and a
nonempty edge list. It is not a cryptographic uniqueness theorem and it does
not assert that the cochain vector is nonzero. For example, the canonical zero
payload is well formed whenever the declared edge count is positive.

Topology compatibility is a separate public invariant. Two committed cochains
can be aggregated only when their moduli, declared edge counts, and topology
hashes agree. The weaker same-topology relation records equality of topology
hashes alone. Cryptographic binding and topology-binding claims are supplied by
the commitment-security axioms, not by the raw payload record.

At the artifact level, Lean records a hashed committed cochain as a public
cochain payload together with digest data and a binding predicate. A
proof-carrying committed cochain adds a proof system, a statement, and a
verification predicate. The basic theorem exposed to this paper is that a
well-formed proof-carrying committed-cochain artifact has an underlying
well-formed cochain commitment payload.

\subsection{Sheaf Storage And Witness Material}

Storage roots, topology hashes, cochain roots, nullifiers, and proof objects are
public artifacts in the stack. Full sheaf, cochain, and fibration data remain
witness material unless a protocol path explicitly exposes them. The storage
mechanics are protocol-level concerns; the mathematical point here is that
stored public artifacts are shadows of accepted witness data, not replacements
for the hidden object.

\subsection{Lossy Projection}

Compression creates verifier-facing public artifacts from hidden or full
witness data. The compression need not be injective; it must preserve enough
structure to support the relevant public statement.

A public shadow is allowed to forget hidden data. In fact, it must forget
enough data that it does not select a private representative. The relevant
question is not whether the public datum is injective, but which facts survive
the compression relation.

\subsection{Preserved And Lost Structure}

The preserved structure before the finite-code layer is:
\begin{enumerate}[label=(\roman*)]
  \item public classifier or fingerprint equality for hidden representatives;
  \item cochain commitment payload well-formedness;
  \item topology binding for cochain commitments;
  \item digest/root binding for compressed public shadows;
  \item proof-carrying validity of the public statement attached to an artifact.
\end{enumerate}

The lost structure is the private representative itself: full sheaf/cochain
witness data, hidden transport, the operative Hodge operator, private
randomness, and runtime lifecycle facts. Those may be bound to public artifacts
or certified by a proof-carrying statement, but they are not made public merely
by the presence of a hash, root, or commitment payload.

\subsection{Fundamental Properties}

\begin{proposition}[Well-formed committed payloads]\label{prop:well-formed-committed-payloads}
  A well-formed cochain commitment payload has nonempty edge data and modulus
  greater than one.
\end{proposition}

\begin{proposition}[Topology-compatible aggregation]\label{prop:topology-compatible-aggregation}
  Commitment aggregation is only meaningful for payloads with compatible
  modulus, declared edge count, and topology hash. This is the public
  topology-binding condition consumed by the chain layer.
\end{proposition}

\begin{proposition}[Proof-carrying artifact exposes payload]\label{prop:proof-carrying-artifact-exposes-payload}
  A well-formed proof-carrying committed-cochain artifact exposes a well-formed
  underlying cochain commitment payload.
\end{proposition}

Byte parsing, allocation bounds, canonical encoding, concrete hash functions,
domain separation, storage inclusion proofs, proof verification, state updates,
and ephemeral-fiber lifecycle properties are TCB obligations. Lean records
their boundary predicates; it does not identify them with the mathematical
cochain commitment payload.

\section{Reed--Solomon Projection And MCA}\label{sec:rs-mca-projection}

The concrete projection chain belongs here because its endpoint is the
Reed--Solomon received-word surface and the mutual-correlated-agreement theorem.
This section uses Grassmannian data in a different role from the HCP
Grassmannian of Sections~\ref{sec:hcp-definition}--\ref{sec:mathematical-framework}.
Earlier, \(W\in\Gr(k,R^n)\) parameterizes hidden harmonic kernels and
projection Laplacians \(\Delta_W=I-P_W\). Here, after the hidden/geometric
object has already been accepted, a frame for the Grassmannian point is used as
the source of a lossy proximity projection through Plucker coordinates and
actual cochain commitments toward a finite Reed--Solomon target.

\subsection{Projection Chain}

The staged chain has the following shape:
\[
\begin{aligned}
  &\text{FrameWitness} \\
  &\quad\to \text{Projective Plucker shadow} \\
  &\quad\to \text{committed cochain/local syndrome} \\
  &\quad\to \text{plane shadow} \\
  &\quad\to \text{finite-field received word} \\
  &\quad\to \text{canonical Reed--Solomon target}.
\end{aligned}
\]
Each stage has local proximity gates and certificates. The orchestrating
projection-chain theorem composes these local certificates without rechecking
them. These are proximity-specific projection gates, not the main \HE{}
acceptance layer from Section~\ref{sec:mathematical-framework}. Their role is
to certify the lossy path into the Reed--Solomon endpoint after the main
hidden/geometric object has already been accepted.
Thus the Plucker shadow is not part of the definition of HCP and is not the
source of hidden-operator hardness. It is the first public projective encoding
used by the proximity proof.

The semantic roles in this chain are separated. The private witness consists of
the fibration data, frame choices, paths, and openings. The shadow is the
Plucker/cochain/proximity data derived from that witness. A hash or root is only
a compressed public handle to canonical bytes for the shadow, not the
mathematical shadow itself. A certificate or SNARK proves that the private
witness produced the accepted shadow and satisfied the stage invariants.

\subsection{Projection-Stage Gates}

The source and Plucker stages carry metric and public-shadow invariants.
We use the usual Grassmannian and Plucker-coordinate algebra: a frame
\((w_1,\ldots,w_r)\) determines the exterior coordinate
\([w_1\wedge\cdots\wedge w_r]\in\mathbb P(\bigwedge^r V)\), with the
quadratic Plucker relations as the projective Grassmannian equations.
The source disagreement score \(\rho(W)\) is required to be nonnegative and
invariant under the public Plucker quotient. The Plucker metric is defined only
on the admissible projective chart, and its distance must be nonnegative and
frame-quotient invariant. Public-shadow respect says that equality of the
public Plucker shadow determines the same committed cochain, plane shadow, and
received word along this projection path.

The Plucker-to-cochain stage carries the local-syndrome interface and consumes
the actual cochain commitment layer of Section~5. Its target is not an abstract
stand-in for a cochain; it is the existing committed payload
\[
  (q,n,c,\tau)
\]
with supported modulus, declared edge count, edge values, and topology hash.
A concrete encoder may send valid Plucker frame data to such a well-formed
payload with the expected edge count, but the storage-native statement is
relational: an observed or proof-carrying cochain commitment \(c\) realizes the
public Plucker shadow. This is the boundary formalized by the public-binding
surface: the chain/storage layer accepts the actual commitment payload, while
the proximity proof supplies the statement that this payload realizes the
Plucker shadow.

Syndrome adequacy then states that a well-formed committed cochain \(c\) with
vanishing local syndrome \(\sigma(c)=0\) is a global-section candidate for the
chosen sheaf. In the concrete bridge, the local syndrome is the
commitment-sum witness; global candidates have zero commitment sum. The point
is that the proximity proof passes through real cochain commitments and their
well-formedness/topology-binding interface, not through a separate synthetic
cochain carrier.

The cochain-to-plane and plane-to-RS stages carry semantic-alignment and
non-collapse conditions. Semantic alignment requires valid frames, the plane
shadow, and the received word to agree with the extraction maps. Plane
extraction records the cochain well-formedness, expected edge count, and
selected-index bounds needed to form the two-coordinate shadow. The far-region
conditions assert calibrated non-collapse from the Plucker source region,
through the committed cochain, to the final Reed--Solomon received word.

\subsection{Canonical Reed--Solomon Endpoint}

The proximity endpoint is stated over the paper-style alphabet shape
\(C \subseteq (\F^s)^n\). The canonical \HE{} evaluation is the ordinary
Reed--Solomon slice at width \(s = 1\), using the standard polynomial
evaluation-code convention.
The canonical endpoint fixes
\[
  \F=\mathbb Z/12289\mathbb Z,\qquad n=4096,\qquad k=1024,
\]
so the rate is \(1/4\). The evaluation domain is a certified smooth power
domain: the chosen generator has order \(12288\), and cubing it gives an
element of order \(4096\). Hence the evaluation points are
\[
  1,\alpha,\alpha^2,\ldots,\alpha^{4095},
\]
with \(\alpha^{4096}=1\) and no earlier repetition. The canonical radius used
for the below-Johnson comparison is the absolute Hamming radius
\[
  2047.
\]
The required finite checks are the field arithmetic, domain cardinality,
smoothness, evaluation-point compatibility, and absolute-radius normalization
for this endpoint; the deep list-decoding and MCA estimates enter only as cited
code-theoretic theorem slots.

\subsection{Mutual Correlated Agreement}

\begin{definition}[MCA bad-challenge fraction]\label{def:mca-bad-challenge}
  Let \(C \subseteq (\F^s)^n\) be a block code and let
  \(\ell(\gamma)=f_1+\gamma f_2\) be an affine line of received words. For a
  radius \(\delta\), the bad challenges are those \(\gamma \in \F\) for which
  \(\ell(\gamma)\) is \(\delta\)-close to \(C\) on a large coordinate set, but
  the pair \((f_1,f_2)\) is not correspondingly close to the two-fold
  interleaved code \(C^{\oplus 2}\) on that same set. The MCA error is the
  fraction of bad challenges in \(\F\).
\end{definition}

For the canonical scalar Reed--Solomon slice \(s=1\), ordinary scalar RS words
are identified with words in \((\F^1)^n\), and the bad-challenge predicate
agrees with the paper-alphabet predicate.

\begin{theorem}[Canonical RS MCA bound]\label{thm:canonical-rs-mca}
  For the canonical \HE{} Reed--Solomon endpoint at the below-Johnson radius,
  every projection-induced final affine line has bad-challenge fraction at most
  \[
    \frac{\tau}{|\F|},
  \]
  where \(\tau\) is the cited MCA threshold for the canonical endpoint.
\end{theorem}

The deep list-decoding and MCA estimates are external theorem inputs: the
Johnson list-decoding bound, the count and fractional MCA bounds, and the paper
Definition~4.3 width-one MCA bound. The internal finite part is the endpoint
plumbing, the \(s=1\) equivalence, and the specialization to projection-induced
lines.
The project-specific list-decoding and MCA input is the cited Proximity Prize
theorem boundary \cite{proximityPrize2026}.

\subsection{Corollary}

\begin{corollary}[Verifier-facing RS/MCA projection]\label{cor:verifier-rs-mca}
  Given an accepted hidden/geometric witness, a certified
  Plucker-to-cochain-to-plane-to-RS projection chain, the required hard
  projection-gate evidence, and a canonical Johnson-radius far-separation
  witness, the public verifier obtains a finite Reed--Solomon received word
  whose canonical MCA bad-challenge fraction is bounded by the cited threshold
  \(\tau/|\F|\).
\end{corollary}

This corollary is a proximity statement. It says that accepted hidden/geometric
state can be projected to a finite public RS/MCA target with calibrated
non-collapse at the selected far region. It does not say that the projection is
injective, and it does not turn the RS/MCA endpoint into the source of HCP
hardness.

\subsection*{Lean Correspondence}

The proximity and finite-code endpoint in this section corresponds to the
following Lean declarations:
\lean{Proximity.GrassmannianPlucker}, \lean{Proximity.CochainBridge},
\lean{Proximity.PluckerCochain}, \lean{Proximity.CochainPlane},
\lean{Proximity.ProjectionChain}, \lean{ProjectionChain.Chain.toEvaluationProjection},
\lean{Proximity.Gates}, \lean{Proximity.GateEvidence},
\lean{CanonicalRS.BridgeData}, \lean{CanonicalRS.CitedBounds},
\lean{CanonicalRS.MCA}, and
\lean{CanonicalRS.MCA.projectionBadChallengeFraction_le}.

\section{Future Directions}

\subsection{Cryptographic Protocol}

The present paper stops at the mathematical and formalization boundary. A
protocol paper should turn the commitment, oracle, proof-carrying, and
proximity surfaces into explicit adversary games and executable verification
paths. The relevant Lean surfaces are already separated: \texttt{TCB} records
proof-carrying artifacts and runtime obligations, \texttt{Gates} records
object-local admissibility predicates, and \texttt{Consensus} records the
Proof-of-Value scoring and policy layer. The concrete protocol must specify
which certificates are public evidence, which witnesses remain private, and
which runtime failures cause rejection, hold, or re-evaluation.

\subsection{Anonymous Identity System}

The hidden-representative/public-classifier split suggests an anonymous
identity architecture. The public identity surface would expose class values,
hashes, and oracle judgments, while the private side would retain fibration
representatives, transport data, and the Hodge operator selecting the operative
kernel. In Lean, this direction is anchored by \texttt{Axioms.HomotopyClass},
\texttt{Axioms.Fibration}, and \texttt{Axioms.OracleModality}. A future system
should make the unlinkability, class-membership, and nullifier properties
explicit rather than treating public-shadow equality as strict identity.

\subsection{Sheaf Storage}

The commitment layer can also be developed as a storage and retrieval system
for controlled sheaf/cochain projections. The public store would retain
topology hashes, cochain roots, compressed shadows, and proof-carrying
statements. Full sheaf and fibration witnesses would remain private or be
opened only through dedicated proof paths. The relevant Lean interfaces are
\texttt{Axioms.CochainCommitment}, \texttt{TCB.CompressionBoundary},
\texttt{TCB.ObservationBoundary}, and \texttt{Proximity.Public}.

\subsection{Vertra Chain}

Vertra is the working name for the chain/network layer. The paper should keep
the mathematical protocol name separate from product and token naming.
The chain layer belongs in the sequel: it consumes the mathematical surfaces
from this paper as protocol evidence, but it also requires block validity,
ledger updates, fork choice, oracle attestation, worker reliability, and
runtime trust-boundary analysis. Lean files such as \texttt{Chain},
\texttt{Consensus}, and \texttt{Consensus.ProofOfValue} already mark this
boundary without making those architecture claims part of the HCP theorem.

\subsection{Complexity-Theoretic Extensions}

The current formalization isolates the hidden-operator obstruction: public data
does not canonically determine the operative Hodge--Laplacian kernel. Future
work should refine this into sharper complexity assumptions, adversary classes,
and oracle-separation statements. Natural targets include quantitative families
of hidden kernels, reductions between HCP variants, leakage models for public
classifiers, and a clearer boundary between external CVP assumptions and
project-native hidden-operator ambiguity.

\appendix

\section{Repository Dependency Graph}

The Lean formalization is organized so that each paper section has a small
number of import surfaces.
\begin{description}[leftmargin=2.4em,style=nextline]
  \item[HCP and CVP.]
    \lean{Axioms/*} feeds \lean{Laplacian/*}, \lean{CVP}, and
    \lean{Hardness}.
  \item[Commitments and public artifacts.]
    \lean{Axioms.CochainCommitment/*} feeds \lean{TCB/*} and
    \lean{Proximity.CochainBridge}.
  \item[Projection chain.]
    \lean{Proximity.GrassmannianPlucker} feeds
    \lean{Proximity.PluckerCochain}, then \lean{Proximity.CochainPlane},
    then \lean{Proximity.CochainReedSolomon}.
  \item[Projection certificates.]
    \lean{Proximity.ProjectionChain} feeds \lean{Proximity.Gates}, which
    feeds \lean{Proximity.CanonicalRS}.
  \item[Protocol boundary.]
    \lean{TCB/*}, \lean{Gates/*}, \lean{Consensus/*}, and
    \lean{Chain} are the sequel-facing protocol surfaces.
\end{description}

For Paper~1, the main Lean check targets are:
\begin{center}
\small
\begin{tabular}{ll}
\lean{lake env lean Axioms.lean} & \lean{lake env lean Gates.lean} \\
\lean{lake env lean TCB.lean} & \lean{lake env lean Hardness.lean} \\
\lean{lake env lean Proximity.lean}
\end{tabular}
\end{center}
The full build target remains \lean{lake build}; the smaller targets above
are the theorem-surface checks most directly corresponding to this paper.

\section{Lean/Paper Correspondence}

\subsection{Definitions And Interfaces}

\begin{description}[leftmargin=2.4em,style=nextline]
  \item[HCP public observation and hidden witness.]
    Sections 2--3; \lean{Hardness.PublicProjection},
    \lean{Hardness.HiddenOperatorFiber}, and
    \lean{Hardness.HiddenOperatorClass}.
  \item[Kernel ambiguity and lossy public data.]
    Sections 2--3; \lean{Hardness.KernelAmbiguous} and
    \lean{Hardness.LossyPublicProjection}.
  \item[Explicit-lattice CVP comparison.]
    Section 3; \lean{CVP.CVP_Instance}, \lean{CVP.CVPSolution}, and
    \lean{CVP.CVPOracleCorrect}.
  \item[Hodge projection model.]
    Sections 2--3; \lean{Laplacian.harmonicProjection},
    \lean{Laplacian.residual}, and \lean{Laplacian.IsHarmonic}.
  \item[Semantic/public equality surface.]
    Section 4; \lean{Axioms.HelferFOL}, \lean{Axioms.HomotopyClass}, and
    \lean{Axioms.Fibration}.
  \item[Heyting/oracle public judgment.]
    Section 4; \lean{Axioms.Heyting3}, \lean{Axioms.OracleQuery},
    \lean{Axioms.OracleAnswer}, and \lean{Axioms.OracleEta2}.
  \item[Cochain commitment payload.]
    Section 5; \lean{Axioms.CochainCommitmentData},
    \lean{CochainCommitmentData.WellFormed}, and
    \lean{CochainCommitmentData.AddCompatible}.
  \item[Commitment security interface.]
    Section 5; \lean{Axioms.CochainCommitmentAxioms}.
  \item[Proof-carrying public artifact.]
    Section 5; \lean{TCB.HashedCommittedCochain} and
    \lean{TCB.ProofCarryingCommittedCochain}.
  \item[Projection chain.]
    Section 6; \lean{ProjectionChain.Chain} and
    \lean{ProjectionChain.Chain.toEvaluationProjection}.
  \item[Proximity-stage certificates.]
    Section 6; \lean{Proximity.Gates.SourceMetricGate},
    \lean{Proximity.Gates.SyndromeAdequacyGate}, and
    \lean{Proximity.Gates.HardGateCertificate}.
  \item[Canonical RS/MCA endpoint.]
    Section 6; \lean{CanonicalRS.BridgeData}, \lean{CanonicalRS.MCA.Line},
    and \lean{CanonicalRS.MCA.badChallengeFraction}.
\end{description}

\subsection{Theorems And Status}

\begin{description}[leftmargin=2.4em,style=nextline]
  \item[Kernel ambiguity implies lossy public projection.]
    \lean{Hardness.kernelAmbiguous_implies_lossy};
    Lean-proved.
  \item[Public-compatible Hodge ambiguity gives kernel ambiguity.]
    \lean{Hardness.kernelAmbiguous_of_publicCompatibleHodgeAmbiguity};
    Lean-proved.
  \item[Hodge solver exposes a known-lattice CVP surface.]
    \lean{Hardness.Hodge_Hardness_Reduction} and
    \lean{Hardness.Hodge_Hardness_Reduction_Proof};
    Lean-proved surface.
  \item[Known-lattice CVP does not select the hidden operator.]
    \lean{Hardness.CVP_Oracle_Insufficient_For_Public_Hodge_Solving};
    Lean-proved surface.
  \item[Projection-model Hodge recovery is CVP-shaped.]
    \lean{Hardness.Hodge_Projection_Is_CVP} and
    \lean{Hardness.harmonicProjection_Is_ProjectionHodgeCVPSolution};
    Lean-proved.
  \item[Exact sheaf gates supply \(d^2=0\) and Laplacian-from-coboundary.]
    \lean{Gates.structural_preconditions_of_sheaf_exact_gates};
    Lean-proved alias/interface.
  \item[Well-formed payloads have nonempty edge data and supported modulus.]
    \lean{CochainCommitmentData.edge_values_nonempty_of_wellFormed} and
    \lean{CochainCommitmentData.modulus_gt_one_of_wellFormed};
    Lean-proved.
  \item[Proof-carrying artifacts expose the committed payload.]
    \lean{TCB.proofCarrying_cochain_wellFormed};
    TCB boundary theorem.
  \item[Projection chain composes to one evaluation projection.]
    \lean{ProjectionChain.Chain.toEvaluationProjection};
    Lean-proved plumbing.
  \item[Syndrome adequacy turns zero local syndrome into codeword/global candidate.]
    \lean{Proximity.Gates.SyndromeAdequacyGate.codeword_of_wellFormed_zeroSyndrome};
    stage-local interface.
  \item[Final proximity evidence yields calibrated end-to-end readiness.]
    \lean{Proximity.Gates.CalibratedFinalCertificate.toCalibratedEndToEndReady};
    Lean-proved plumbing.
  \item[Canonical RS list/MCA bounds enter at a cited boundary.]
    \lean{CanonicalRS.CitedBounds.bounds};
    cited external theorem package.
  \item[Canonical MCA bad-challenge fraction is bounded after projection.]
    \lean{CanonicalRS.MCA.projectionBadChallengeFraction_le};
    Lean-proved from cited bounds.
\end{description}


\bibliographystyle{alpha}
\bibliography{references}

\end{document}
